% Options for packages loaded elsewhere
\PassOptionsToPackage{unicode}{hyperref}
\PassOptionsToPackage{hyphens}{url}
\PassOptionsToPackage{dvipsnames,svgnames,x11names}{xcolor}
\documentclass[
  11pt,
]{article}
\usepackage{xcolor}
\usepackage[a4paper,left=24mm,right=24mm,top=24mm,bottom=24mm]{geometry}
\usepackage{amsmath,amssymb}
\setcounter{secnumdepth}{-\maxdimen} % remove section numbering
\usepackage{iftex}
\ifPDFTeX
  \usepackage[T1]{fontenc}
  \usepackage[utf8]{inputenc}
  \usepackage{textcomp} % provide euro and other symbols
\else % if luatex or xetex
  \usepackage{unicode-math} % this also loads fontspec
  \defaultfontfeatures{Scale=MatchLowercase}
  \defaultfontfeatures[\rmfamily]{Ligatures=TeX,Scale=1}
\fi
\usepackage{lmodern}
\ifPDFTeX\else
  % xetex/luatex font selection
  \setmainfont[]{Times New Roman}
  \setsansfont[]{Arial}
  \setmonofont[]{Menlo}
\fi
% Use upquote if available, for straight quotes in verbatim environments
\IfFileExists{upquote.sty}{\usepackage{upquote}}{}
\IfFileExists{microtype.sty}{% use microtype if available
  \usepackage[]{microtype}
  \UseMicrotypeSet[protrusion]{basicmath} % disable protrusion for tt fonts
}{}
\makeatletter
\@ifundefined{KOMAClassName}{% if non-KOMA class
  \IfFileExists{parskip.sty}{%
    \usepackage{parskip}
  }{% else
    \setlength{\parindent}{0pt}
    \setlength{\parskip}{6pt plus 2pt minus 1pt}}
}{% if KOMA class
  \KOMAoptions{parskip=half}}
\makeatother
\usepackage{longtable,booktabs,array}
\usepackage{caption}
\captionsetup[table]{skip=6pt}
\newcounter{none} % for unnumbered tables
\usepackage{calc} % for calculating minipage widths
% Correct order of tables after \paragraph or \subparagraph
\usepackage{etoolbox}
\makeatletter
\patchcmd\longtable{\par}{\if@noskipsec\mbox{}\fi\par}{}{}
\makeatother
% Allow footnotes in longtable head/foot
\IfFileExists{footnotehyper.sty}{\usepackage{footnotehyper}}{\usepackage{footnote}}
\makesavenoteenv{longtable}
\usepackage{graphicx}
\makeatletter
\newsavebox\pandoc@box
\newcommand*\pandocbounded[1]{% scales image to fit in text height/width
  \sbox\pandoc@box{#1}%
  \Gscale@div\@tempa{\textheight}{\dimexpr\ht\pandoc@box+\dp\pandoc@box\relax}%
  \Gscale@div\@tempb{\linewidth}{\wd\pandoc@box}%
  \ifdim\@tempb\p@<\@tempa\p@\let\@tempa\@tempb\fi% select the smaller of both
  \ifdim\@tempa\p@<\p@\scalebox{\@tempa}{\usebox\pandoc@box}%
  \else\usebox{\pandoc@box}%
  \fi%
}
% Set default figure placement to htbp
\def\fps@figure{htbp}
\makeatother
\setlength{\emergencystretch}{3em} % prevent overfull lines
\providecommand{\tightlist}{%
  \setlength{\itemsep}{0pt}\setlength{\parskip}{0pt}}
\usepackage{fancyhdr}
\usepackage{titlesec}
\usepackage{float}
\usepackage{needspace}
\definecolor{ink}{HTML}{163B4A}
\definecolor{muted}{HTML}{54717B}
\pagestyle{fancy}
\fancyhf{}
\fancyhead[L]{\sffamily\scriptsize\color{muted} VISUAL ATTENTION AND WORKING MEMORY}
\fancyhead[R]{\sffamily\scriptsize\color{muted} TEXTBOOK COMPANION}
\fancyfoot[R]{\sffamily\small\thepage}
\renewcommand{\headrulewidth}{0.3pt}
\setlength{\headheight}{15pt}
\titleformat{\section}{\sffamily\Large\bfseries\color{ink}}{}{0pt}{}
\titleformat{\subsection}{\sffamily\large\bfseries\color{ink}}{}{0pt}{}
\titlespacing*{\section}{0pt}{8pt}{14pt}
\titlespacing*{\subsection}{0pt}{15pt}{7pt}
\setlength{\parskip}{6pt}
\setlength{\parindent}{0pt}
\setlength{\emergencystretch}{3em}
\setlength{\tabcolsep}{5pt}
\renewcommand{\arraystretch}{1.15}
\floatplacement{figure}{H}
\widowpenalty=10000
\clubpenalty=10000
\AtBeginEnvironment{longtable}{\small}
\let\oldsubsection\subsection
\renewcommand{\subsection}{\Needspace{6\baselineskip}\oldsubsection}
\AtBeginDocument{\hypersetup{colorlinks=true,linkcolor=ink,urlcolor=ink}}
\usepackage{bookmark}
\IfFileExists{xurl.sty}{\usepackage{xurl}}{} % add URL line breaks if available
\urlstyle{same}
% fallback for those not using the hyperref driver hyperxmp:
\makeatletter
\@ifundefined{xmpquote}{\newcommand{\xmpquote}[1]{#1}}{}
\makeatother
\hypersetup{
  pdftitle={Visual attention and working memory},
  pdfauthor={Project reading edition},
  colorlinks=true,
  linkcolor={ink},
  filecolor={Maroon},
  citecolor={Blue},
  urlcolor={ink},
  pdfcreator={LaTeX via pandoc}}

\title{Visual attention and working memory}
\usepackage{etoolbox}
\makeatletter
\providecommand{\subtitle}[1]{% add subtitle to \maketitle
  \apptocmd{\@title}{\par {\large #1 \par}}{}{}
}
\makeatother
\subtitle{A textbook companion to our KDA models, spatial recurrence,
and training tasks}
\author{Project reading edition}
\date{24 September 2026}

\begin{document}
\maketitle

{
\setcounter{tocdepth}{2}
\tableofcontents
}
\section{How to use this reading}\label{how-to-use-this-reading}

This is a teaching companion to \emph{Spatial KDA memory: architecture
and causal tests}, the ten-page note dated 23 September 2026. That note
is organized around reporting and experimental recommendations. This
book is organized around understanding: what computation each component
performs, why it is there, what a complete trial asks the model to do,
and what we can infer from its answer.

You are assumed to know neural-network training, convolution, matrix
multiplication and probability. We will not rehearse those basics. We
will unpack the less familiar machinery: associative matrix memory,
error-correcting writes, channel-wise retention, spatial recurrence,
recurrent readout, and the exact semantics of the task suite. Equations
are followed by their operational meaning and, where useful, a small
worked example.

\textbf{Architecture status matters.} The existing joint learner uses
three spatial KDA modules followed by a vector GRU. The new branch being
built by another agent retains those KDA modules and replaces the final
vector-recurrent route with a spatial ConvGRU. This book describes that
branch as a \textbf{planned design}, according to its recorded brief and
your clarification. It makes no claim that the new branch is
implemented, training, or better. Historical ConvGRU accumulator
experiments are a third, different architecture.

The source PDF's behavioral numbers are a dated snapshot, not a live
progress report. We preserve its scientific questions but do not turn
its interim scores into current results. No training, model evaluation,
checkpoint modification or intervention was performed to write this
reading. Illustrative numbers below are calculations or constructed
examples, explicitly distinguished from measured model behavior.

Read Chapters 1-5 consecutively for the architecture. Chapters 6-7 teach
all thirteen tasks. Chapter 8 connects the tasks to learning and
evaluation. Chapter 9 returns to the causal questions in the original
note. Chapter 10 provides exercises and worked answers. The source map
at the end identifies the implementation or design behind each account.

\newpage

\section{1. The computation a trial
requires}\label{the-computation-a-trial-requires}

\subsection{1.1 From pixels to an
answer}\label{from-pixels-to-an-answer}

Consider a delayed orientation trial. A glyph indicates a location and a
requested rotation sign. Four oriented patches appear. They disappear
for several frames, then reappear. The required answer is whether the
selected patch rotated in the requested sign.

A useful decomposition is \textbf{encode, select, retain, compare,
report}. Encoding must represent orientation well enough to distinguish
the relevant rotation. Selection must associate the instruction with the
correct spatial patch. Retention must preserve useful sample information
while the patch is absent. Comparison must relate that information to
the probe. Reporting must map the comparison and instruction to the
appropriate class.

Those are functional requirements, not five named modules in our code. A
recurrent network can distribute one function across several states or
use one state for several functions. For example, the final GRU might
retain the cue, while KDA retains local appearance. Alternatively, the
GRU might retain a partly computed answer. The architecture permits
these possibilities; behavioral accuracy alone does not tell us which
one was learned.

This decomposition prevents an easy mistake: calling every failure after
a delay a memory failure. If the network never represented orientation
accurately, perfect retention would preserve an inadequate
representation. If it retained the right information but used the wrong
location, the failure would look similar at the output. We need tasks
and diagnostics that separate these explanations.

\subsection{1.2 Three kinds of
persistence}\label{three-kinds-of-persistence}

\textbf{Learned parameters} persist between trials. They include
convolution kernels, the networks that generate KDA keys and gates,
recurrent weights and task-head weights. During training, the optimizer
changes them across batches. At inference, they define the rules of
computation.

\textbf{Recurrent activation state} persists within a trial. KDA's
matrices and the GRU hidden state change as frames arrive. They are
reset between trials. They are not optimizer parameters. A matrix that
changes at every frame is therefore not evidence that the model is
changing its trained weights at inference.

\textbf{The raw-frame stack} makes the previous two images directly
available with the current image. This is a fixed buffer, not a learned
retention mechanism. Its influence must be separated from persistence in
KDA or GRU state. {[}C1-C3{]}

A familiar associative-memory interpretation calls a rapidly updated
matrix a \emph{fast weight}. That term is useful only if we remember the
distinction: here it is a trial-local activation computed by a learned
rule, not an additional optimizer-trained parameter tensor.

\subsection{1.3 What the model receives}\label{what-the-model-receives}

The suite produces an image tensor of shape

\[X\in[0,1]^{B\times T\times3\times100\times100}.\]

Here \(B\) counts independent trials and \(T\) counts presented frames.
A batch has one task and one condition, so all trials have the same
native length. The model centers pixels by subtracting \(0.5\) and
constructs

\[x_t=\operatorname{concat}(X_{t-2}-0.5,\;X_{t-1}-0.5,\;X_t-0.5).\]

Missing earlier frames are zero in centered coordinates. The encoder
therefore receives nine channels at each step. This operation is causal:
it never includes a future frame. On a two-frame task, there are still
only two actual observations. Padding is not an extra stimulus.
{[}C2-C4{]}

Only images enter the shared visual and recurrent computation. The
experimenter selects the output head using the task ID. Numerical
angles, target indices, phase labels, image identities and correct
answers are analysis metadata, not input features. Visible instruction
glyphs, however, are pixels and can be learned from. The model is not
required to infer the externally selected task ID.

\subsection{1.4 A blank frame does not immediately remove the
sample}\label{a-blank-frame-does-not-immediately-remove-the-sample}

Suppose \(s\) is the final sample frame and \(b_1,b_2,b_3\) are
subsequent blanks. The first blank update still sees two earlier images;
the second sees \([s,b_1,b_2]\); only the third sees \([b_1,b_2,b_3]\).
Thus the first two nominal blanks still offer direct sample access
through the stack.

For the ring task, the final probe step sees both sample frames and the
probe together. That task can test spatially selective comparison
without establishing learned retention through a long absence. At longer
delays, the sample leaves the stack and useful information must reach
the answer through other routes.

Blank processing is also not an identity operation. Convolution biases,
instruction markers, normalization, KDA writes and recurrent dynamics
can remain active. A model must learn appropriate behavior on blanks; it
does not automatically freeze when the stimulus disappears.

\textbf{Check your understanding.} If a model is perfect at zero delay
but fails after four blank frames, what has been established? There is a
dependence on the longer trial or stimulus absence. We have not yet
distinguished failed retention from changed state dynamics, cue loss, or
an unfamiliar comparison context.

\newpage

\section{2. KDA: learning how to revise an associative
memory}\label{kda-learning-how-to-revise-an-associative-memory}

\subsection{2.1 The lineage and the local
implementation}\label{the-lineage-and-the-local-implementation}

KDA means \textbf{Kimi Delta Attention}. The Kimi Linear paper develops
a gated delta-rule memory with fine-grained decay, within a larger
language-model architecture. Our visual model adopts a small spatial
version of the associative recurrence. It does not reproduce the full
Kimi Linear architecture, its hybrid attention layers, or its optimized
execution system. The equations below follow our implementation; the
derivations and numerical examples explain that code. {[}C1; P1{]}

\subsection{2.2 A matrix represents
associations}\label{a-matrix-represents-associations}

At one site and one head, let the state be
\(S\in\mathbb R^{d_k\times d_v}\). A key \(k\in\mathbb R^{d_k}\)
addresses an association, while a value \(v\in\mathbb R^{d_v}\) is the
content associated with that key. Given a query \(q\in\mathbb R^{d_k}\),
the read is

\[o=S^\top q\in\mathbb R^{d_v}.\]

Why does this act like memory? Imagine first writing one outer product
\(S=kv^\top\). Then

\[S^\top q=v(k^\top q).\]

The query returns the value multiplied by its alignment with the key. If
the normalized query equals the normalized key, it returns \(v\). If the
query is orthogonal to the key, this association contributes zero.
Queries and keys are learned features, so their coordinates need not
correspond to interpretable things such as ``red,'' ``left,'' or
``orientation.''

For several additive writes,

\[S=\sum_i k_i v_i^\top,\qquad S^\top q=\sum_i(k_i^\top q)v_i.\]

This is an associative retrieval rule. It resembles attention in that a
query controls the contribution of stored content. These coefficients
need not be positive or sum to one. They are not probabilities over
locations, and storing this matrix does not preserve a separate record
of every past image.

\subsection{2.3 Why plain addition is
insufficient}\label{why-plain-addition-is-insufficient}

Suppose the same unit key repeatedly appears with the same value. The
additive rule \(S_t=S_{t-1}+k_tv_t^\top\) accumulates that value again
each time. After ten identical writes from zero, querying that key
returns ten times the original value. Sometimes accumulation is wanted;
sometimes repeated observation should confirm a stable association
without making it grow indefinitely.

The delta rule asks a more selective question: \textbf{What does memory
currently predict at this key, and what correction is needed?} It writes
the residual error rather than the whole observation. A repeated
association that is already predicted well needs little additional
correction.

This distinction is relevant to our tasks. Repeated sample frames need
not be separate items. Repeated recognition probes should not
automatically count as repeated study evidence. An error-based write
creates the possibility of controlled revision, although the network
must still learn representations and gates that use that possibility
effectively.

\subsection{2.4 Deriving the error-correcting
write}\label{deriving-the-error-correcting-write}

For the moment, ignore decay. Define an instantaneous reconstruction
objective

\[\ell(S)=\tfrac12\|v-S^\top k\|_2^2.\]

Its gradient with respect to \(S\) is

\[\nabla_S\ell=k(S^\top k-v)^\top.\]

A step of size \(\beta\) in the negative-gradient direction gives

\[S_{\mathrm{new}}=S+\beta k(v-S^\top k)^\top.\]

That is the delta-rule write. This derivation explains the form of the
recurrence; our training loop does \textbf{not} separately minimize this
local objective with another optimizer. The code directly computes the
update. The global task loss trains the networks that produce \(k,v,q\)
and the gates.

To see the correction precisely, set \(a=S^\top k\). After the write,

\[S_{\mathrm{new}}^\top k=a+\beta\|k\|_2^2(v-a).\]

For a unit key this becomes \((1-\beta)a+\beta v\). The retrieved
prediction moves a fraction \(\beta\) toward the new value. With
\(\beta=1\), it exactly fits that key's current value. With a very small
\(\beta\), it changes little. The implementation uses sigmoid gates, so
exact zero and one are limiting cases rather than ordinary finite-logit
outputs.

\subsection{2.5 Retention comes before the
correction}\label{retention-comes-before-the-correction}

Our full recurrence has four operations:

\[\bar S_t=D_{\alpha_t}S_{t-1},\qquad
\hat v_t=\bar S_t^\top k_t,\] \[e_t=v_t-\hat v_t,\qquad
S_t=\bar S_t+\beta_t k_t e_t^\top,\] \[o_t=S_t^\top q_t.\]

\(D_{\alpha_t}\) is diagonal. Each of its \(d_k\) entries scales one row
of the old matrix. The memory first retains or attenuates old
associations, then predicts from that decayed state, then corrects the
prediction, then reads the updated state. Reading after writing means
the current observation can immediately influence the emission.

The two gates have different jobs. \textbf{Retention \(\alpha\)}
controls survival of existing state along key coordinates. \textbf{Write
strength \(\beta\)} controls the size of the current error correction. A
small write gate does not preserve memory if retention is also small. A
large retention gate does not prevent a strong write from overwriting an
overlapping association.

Retention is row-wise, not one independently controlled decay per value
element. In our model there are eight retention values per head, and
each scales a row containing sixteen value coordinates. The learned key
basis determines what those rows mean.

\begin{figure}
\centering
\includegraphics[width=0.95\linewidth,height=\textheight,keepaspectratio,alt={The associative update separates retention, prediction, correction and readout. Each operation has a different role.}]{tmp/pdfs/kda_step.pdf}
\caption{The associative update separates retention, prediction,
correction and readout. Each operation has a different role.}
\end{figure}

\subsection{2.6 A complete numerical
step}\label{a-complete-numerical-step}

Use a two-key, two-value toy memory, smaller than the deployed module:

\[S_{t-1}=\begin{bmatrix}2&0\\0&3\end{bmatrix},\quad
\alpha=(0.5,1),\quad k=q=\begin{bmatrix}1\\0\end{bmatrix},\quad
v=\begin{bmatrix}4\\2\end{bmatrix},\quad\beta=0.5.\]

Decay gives
\(\bar S_t=\left[\begin{smallmatrix}1&0\\0&3\end{smallmatrix}\right]\).
The current key predicts \((1,0)^\top\), so the error is \((3,2)^\top\).
The rank-one correction is

\[\beta ke^\top=\begin{bmatrix}1.5&1\\0&0\end{bmatrix}.\]

Therefore

\[S_t=\begin{bmatrix}2.5&1\\0&3\end{bmatrix},\qquad o_t=(2.5,1)^\top.\]

The output is halfway between the decayed prediction \((1,0)\) and the
requested value \((4,2)\). The second row is unaffected by this write
because the key points entirely along the first coordinate. This is a
constructed arithmetic example, not an extracted model state.

\subsection{2.7 Interference follows from key
overlap}\label{interference-follows-from-key-overlap}

Now ask what this write does to another query \(q'\). Its immediate
change relative to the decayed state is

\[\Delta o(q')=\beta(k^\top q')e.\]

Orthogonal queries are unaffected; aligned queries receive the full
correction; negatively aligned queries receive a signed effect. There is
no guarantee that one object occupies one orthogonal address. Learned
keys can overlap, and a write for one stimulus can alter retrieval for
another.

The matrix's rank is at most \(\min(d_k,d_v)\). Our \(8\times16\) matrix
has rank at most eight. That does \textbf{not} mean it stores exactly
eight objects. Object capacity depends on the representational code,
interference, required precision, downstream readout and task
distribution. ``Eight key dimensions'' is an architectural fact; ``eight
remembered items'' would be an empirical claim requiring a capacity
experiment.

\subsection{2.8 Retention is not a single fixed time
constant}\label{retention-is-not-a-single-fixed-time-constant}

If writes were absent and every row retained the same fixed fraction
\(a\), old state would decay as \(a^n\) after \(n\) steps. For
illustration, \(0.9^{24}\approx0.080\) whereas
\(0.99^{24}\approx0.786\). Small differences in retention can strongly
affect long delays.

Our retention and write gates are input-dependent after training, and
different rows can have different values. Blank frames can generate
nonzero values and writes. Successive state changes therefore cannot
generally be summarized by one exponential. Initial retention is \(0.9\)
and initial write strength is \(0.5\), but these are initialization
settings, not measured trained laws. {[}C1{]}

\subsection{2.9 What the stability argument actually
proves}\label{what-the-stability-argument-actually-proves}

Expanding the correction gives

\[S_t=A_tS_{t-1}+\beta_t k_tv_t^\top,\qquad
A_t=(I-\beta_tk_tk_t^\top)D_{\alpha_t}.\]

For fixed keys and gates, with \(\|k_t\|_2\leq1\) and
\(0\leq\beta_t\leq1\), the rank-one factor has eigenvalue
\(1-\beta_t\|k_t\|^2\) along the key and eigenvalue one in orthogonal
directions. Its spectral norm is at most one. Hence

\[\|A_t\|_2\leq\max_j\alpha_{t,j}\leq1.\]

If two memories receive the \emph{same fixed input-derived quantities},
their difference contracts or stays unchanged under this homogeneous
transition. This is useful, but limited. It is not a proof that the
complete network's gradients are bounded. Values are unconstrained; the
input-driven term can add state; and the hierarchy contains other
memory-dependent computations and a final recurrent route. Nor does
nonexpansion guarantee accurate memory: a map that aggressively erases
information is very stable.

\subsection{2.10 The temporal attention hidden inside the
recurrence}\label{the-temporal-attention-hidden-inside-the-recurrence}

With \(S_0=0\), unrolling the recurrence yields

\[S_t=\sum_{i=1}^{t}(A_tA_{t-1}\cdots A_{i+1})\beta_i k_iv_i^\top.\]

An empty product is the identity. Querying gives

\[o_t=\sum_{i=1}^{t}w_{t,i}v_i,\qquad
w_{t,i}=\beta_iq_t^\top(A_t\cdots A_{i+1})k_i.\]

Each past value contributes according to its write strength, its
surviving address after intervening transitions, and the present query.
These are \textbf{implicit temporal contributions} at one site/head.
They need not sum to one, can be negative, and are not a softmax
allocation over the image. Moreover, changing an actual past image
changes downstream keys and gates too; this fixed-trajectory expansion
is not by itself a causal attribution to pixels.

\newpage

\section{3. Spatial KDA: a field of small associative
memories}\label{spatial-kda-a-field-of-small-associative-memories}

\subsection{3.1 What makes it spatial}\label{what-makes-it-spatial}

A sequence KDA module maintains a matrix as tokens arrive. Our spatial
adaptation maintains a separate matrix at every location of a feature
map, updating those matrices as video frames arrive. Time is the
recurrent axis; image positions remain explicit indices.

At scale \(s\), location \(p\) and head \(h\),

\[S^{(s)}_{t,p,h}\in\mathbb R^{8\times16},\qquad h\in\{1,2\}.\]

The full tensor is \(B\times H_s\times W_s\times2\times8\times16\). The
two heads are separate learned associative systems at the same site.
They are not left/right stimulus slots or experimenter-assigned feature
types. Both follow the same algebra but have different learned
projections. {[}C1-C2{]}

A site stores 256 scalars, but emits 32: two sixteen-dimensional queried
values. Thus the downstream network does not receive the whole state
matrix. It receives a query-dependent view of it. Information might
remain in the state while the current query fails to expose it.

\subsection{3.2 How pixels become keys and
gates}\label{how-pixels-become-keys-and-gates}

The encoder block produces \(C_s\) channels. A learned \(1\times1\)
projection maps these to 32 channels at each location. A \(3\times3\)
convolution then generates 82 channels, packed as two groups of 41:

\[41=8\;(q)+8\;(k)+16\;(v)+8\;(\alpha)+1\;(\beta).\]

Queries and keys are L2-normalized, using epsilon \(10^{-6}\) to handle
very small norms. The gates use sigmoids. Values have no such bound.
After the update and read, the two head outputs are concatenated and
transformed by a learned \(1\times1\) convolution, including a bias.
This yields the emitted field \(O_t\). {[}C1{]}

This packing explains what ``input-dependent gates'' means concretely:
changing a local input neighborhood changes convolution outputs, which
can change retention and write strength. The same-scale previous matrix
is not concatenated into that gate-producing convolution. It influences
prediction and readout through the recurrence. At higher scales, current
inputs already include lower-scale memory emissions, so higher-scale
gating can be memory-conditioned indirectly.

\subsection{3.3 Local state does not imply an isolated
pixel}\label{local-state-does-not-imply-an-isolated-pixel}

There is no direct term transporting \(S_{t-1,p'}\) into \(S_{t,p}\)
from a neighboring site in the KDA update itself. Nevertheless, the
quantities that control that update come from spatial convolutions. Each
site represents a receptive field, not one raw pixel. Deeper scales
combine broader neighborhoods, and GroupNorm uses spatial statistics
within channel groups.

These facts matter for interpreting a map. A high gate value near a
target is not automatically a neuron specifically coding that target. A
perturbation applied to a few feature sites can affect downstream sites,
normalization statistics and the final decision. ``Spatially indexed''
describes how the state is organized; it does not certify strict causal
independence among regions.

\subsection{3.4 State, emission and decision are different
objects}\label{state-emission-and-decision-are-different-objects}

The \textbf{state} is the matrix carried to the next frame. The
\textbf{raw read} is \(S_t^\top q_t\). The \textbf{emission} applies the
learned output projection to the concatenated reads. The
\textbf{decision representation} is computed after further convolution
and recurrence.

These distinctions explain several otherwise puzzling observations. Zero
state does not necessarily imply zero emission, because the output
projection has a bias. A large state norm can produce a small read if
the query is nearly orthogonal to its important directions. A visible
emission change can be ignored by downstream weights. Conversely, a
small targeted feature change can have a large decision effect near a
classification boundary.

None of these quantities is a universal ``attention strength.'' Gate
magnitude, state norm, query-key alignment, feature influence and
behavioral selection describe different aspects of computation. A
textbook account should make it possible to say which one is being
plotted.

\subsection{3.5 Storage and computation
costs}\label{storage-and-computation-costs}

The three KDA grids contain \(25^2+13^2+7^2=843\) sites. Multiplying by
256 gives \textbf{215,808 state scalars per example}, or 863,232 bytes
in fp32. This excludes frame buffers, the final recurrent state,
intermediate activations and optimizer storage.

Each KDA module has \(32\times82\times3\times3+82=23,698\)
input-convolution parameters and \(32\times32+32=1,056\)
output-projection parameters: \textbf{24,754 learned parameters per
scale}. The input projection from \(C_s\) to 32 is separate. These
arithmetic counts follow the source shapes. {[}C1-C2{]}

At fixed image resolution, recurrent state size does not grow with
sequence length. Full backpropagation through a longer trial still needs
more training-time intermediates. Constant-size forward state therefore
does not mean constant training memory or constant work for an entire
sequence.

\newpage

\section{4. The implemented joint
architecture}\label{the-implemented-joint-architecture}

\subsection{4.1 The visual hierarchy, one frame at a
time}\label{the-visual-hierarchy-one-frame-at-a-time}

The existing joint learner uses the following route. ``Visual'' is the
encoder output before concatenation; ``combined'' is what the next stage
receives. Every learned component is trainable. {[}C2-C3{]}

{\def\LTcaptype{none} % do not increment counter
\begin{longtable}[]{@{}
  >{\raggedright\arraybackslash}p{(\linewidth - 6\tabcolsep) * \real{0.2500}}
  >{\raggedright\arraybackslash}p{(\linewidth - 6\tabcolsep) * \real{0.2500}}
  >{\raggedright\arraybackslash}p{(\linewidth - 6\tabcolsep) * \real{0.2500}}
  >{\raggedright\arraybackslash}p{(\linewidth - 6\tabcolsep) * \real{0.2500}}@{}}
\toprule\noalign{}
\begin{minipage}[b]{\linewidth}\raggedright
Stage
\end{minipage} & \begin{minipage}[b]{\linewidth}\raggedright
Visual output
\end{minipage} & \begin{minipage}[b]{\linewidth}\raggedright
Memory emission
\end{minipage} & \begin{minipage}[b]{\linewidth}\raggedright
Combined output
\end{minipage} \\
\midrule\noalign{}
\endhead
\bottomrule\noalign{}
\endlastfoot
Conv 1 & \(32\times50\times50\) & None & \(32\times50\times50\) \\
Conv 2 & \(64\times25\times25\) & \(32\times25\times25\) &
\(96\times25\times25\) \\
Conv 3 & \(96\times13\times13\) & \(32\times13\times13\) &
\(128\times13\times13\) \\
Conv 4 & \(128\times7\times7\) & \(32\times7\times7\) &
\(160\times7\times7\) \\
\end{longtable}
}

All convolutions have stride two. The first kernel is \(5\times5\); the
others are \(3\times3\), with padding. Each block uses eight-group
GroupNorm followed by ReLU. GroupNorm here is an encoder operation; it
does not directly normalize every KDA state matrix.

The order is important. At a given frame, Conv 2 is followed by the
first KDA update and read. That emission becomes part of the input to
Conv 3, whose output drives the next KDA. Processing then repeats at
Conv 4. Memory therefore participates inside the hierarchy, rather than
being appended only after all visual computation is finished.

\subsection{4.2 The final compression}\label{the-final-compression}

Let \(F_t\in\mathbb R^{160\times7\times7}\) be the deepest combined map.
The old route is

\[f_t=\operatorname{ReLU}(W_f\operatorname{vec}(F_t)+b_f)\in\mathbb R^{256},\]
\[h_t=\operatorname{GRU}(f_t,h_{t-1})\in\mathbb R^{256},\qquad
\ell=W_{\mathrm{task}}h_T+b_{\mathrm{task}}.\]

The linear projection compresses 7,840 map coordinates to 256 features
\textbf{before} the final recurrent update. Flattening does not remove
location labels in the way global average pooling would: different
positions get different learned weights. But the final recurrent state
has no explicit spatial grid, and each time step's map must pass through
that compressed representation.

This motivates a readout hypothesis. Perhaps the spatial memories
contain useful information, but the final route does not preserve or
compare it adequately. It remains a hypothesis. Compression can also
help by extracting sufficient task statistics and removing irrelevant
detail. A dimensional bottleneck alone is not evidence that the model
fails because of that bottleneck.

\subsection{4.3 The vector GRU's gates}\label{the-vector-grus-gates}

A GRU combines a previous state with a candidate. In the PyTorch
convention used by the existing \texttt{nn.GRU}, the update gate \(z_t\)
weights the \textbf{old} state:

\[r_t=\sigma(W_{ir}f_t+b_{ir}+W_{hr}h_{t-1}+b_{hr}),\]
\[z_t=\sigma(W_{iz}f_t+b_{iz}+W_{hz}h_{t-1}+b_{hz}),\]
\[n_t=\tanh\!\left(W_{in}f_t+b_{in}+r_t\odot(W_{hn}h_{t-1}+b_{hn})\right),\]
\[h_t=(1-z_t)\odot n_t+z_t\odot h_{t-1}.\]

The reset gate controls the recurrent contribution to the candidate; the
update gate controls mixing with the old state. PyTorch applies the
reset multiplication after the hidden affine transform in the candidate.
This detail differs from placing the reset on the hidden state before
the transform. {[}P3{]}

The proposed ConvGRU will use a \textbf{write} gate instead, where a
larger value weights the candidate more. When comparing plots or
equations, compare gate roles, not names. A large old-state update gate
and a large candidate-write gate imply opposite mixing tendencies.

\subsection{4.4 Why more than one memory route
matters}\label{why-more-than-one-memory-route-matters}

The model can retain information in three KDA fields, the vector GRU,
and the two-frame input history. It has no explicit top-down connection
from final GRU state back to the encoder. The KDA states nevertheless
affect subsequent visual processing through their emissions.

Suppose resetting all KDA matrices worsens accuracy. That establishes an
effect of those states under that intervention, but it does not mean all
useful memory resided exclusively in KDA. The GRU may depend on KDA's
normal outputs and react badly when they are changed. Conversely, an
apparently harmless reset may be masked by the raw-frame stack or
redundancy in the GRU.

The final task head is linear. Complex spatial comparison must therefore
already be encoded in \(h_T\) in a form the head can separate. There is
no explicit hand-programmed comparator, object-slot matching mechanism,
or experimenter-supplied target coordinate in this route.

\subsection{4.5 Reading the parameter count
correctly}\label{reading-the-parameter-count-correctly}

The existing thirteen-head model has 2,750,324 trainable parameters:
256,992 in encoder blocks, 9,312 in KDA input projections, 74,262 in KDA
modules, 2,007,296 in the dense feature projection, 394,752 in the
vector GRU and 7,710 in task heads. These are static architecture
counts, confirmed in the original report's source audit. {[}C1-C3; R1{]}

Most parameters are in the dense feature projection, while most
trial-local recurrent scalars are in KDA. Parameter count measures
learned degrees of freedom in the computation; activation-state size
measures the amount of intermediate state carried for a particular
example. They answer different questions.

\newpage

\section{5. ConvGRU and the planned spatial
readout}\label{convgru-and-the-planned-spatial-readout}

\subsection{5.1 What is actually being
changed}\label{what-is-actually-being-changed}

\textbf{Design status: planned, under construction by another agent.}
The recorded proposal retains the centered three-frame stack,
convolutional hierarchy, all three KDA modules and thirteen task heads.
It replaces the final dense projection plus vector GRU with a spatial
projection and a ConvGRU. The final dense projection moves to after
recurrent processing. {[}D1{]}

The proposed route is

\[F_t\in\mathbb R^{160\times7\times7}
\xrightarrow{1\times1\;\mathrm{conv}}U_t\in\mathbb R^{64\times7\times7},\]
\[H_t=\operatorname{ConvGRU}(U_t,H_{t-1}),\quad
H_t\in\mathbb R^{64\times7\times7},\]
\[a_T=\operatorname{ReLU}(W_a\operatorname{vec}(H_T)+b_a)\in\mathbb R^{256},\qquad
\ell=W_{\mathrm{task}}a_T+b_{\mathrm{task}}.\]

The essential change is \textbf{where spatial information is
compressed}. The old route projects the deepest map to a 256-vector at
every frame and then recurs. The proposed route recurs on a \(7\times7\)
grid and projects only the final spatial state. This provides the final
recurrent computation with an explicit spatial organization throughout
the trial.

It does not abolish compression. The \(1\times1\) projection reduces 160
channels to 64, the recurrent field has finite capacity, and the final
head still receives only 256 features. It also does not add an explicit
target-selection controller or top-down feedback.

\begin{figure}
\centering
\includegraphics[width=0.9\linewidth,height=\textheight,keepaspectratio,alt={The planned branch moves the final spatial compression to after recurrent processing. It retains the same KDA hierarchy.}]{tmp/pdfs/architecture_routes.pdf}
\caption{The planned branch moves the final spatial compression to after
recurrent processing. It retains the same KDA hierarchy.}
\end{figure}

\subsection{5.2 The ConvGRU equations}\label{the-convgru-equations}

Convolutional recurrence uses shared spatial kernels to combine a
current map and an earlier hidden map. Ballas and colleagues provide a
primary precedent for ConvGRU video representations; our exact gate
convention and dimensions come from the local design brief. {[}P2; D1{]}

Let \(*\) denote convolution and \([U_t,H_{t-1}]\) concatenation along
channels. The planned cell computes

\[[w_t,r_t]=\sigma\!\left(K_g*[U_t,H_{t-1}]+b_g\right),\]
\[\widetilde H_t=\tanh\!\left(K_c*[U_t,r_t\odot H_{t-1}]+b_c\right),\]
\[H_t=(1-w_t)\odot H_{t-1}+w_t\odot\widetilde H_t.\]

The gate convolution has 128 input channels and 128 output channels,
split into 64 write gates and 64 reset gates. The candidate convolution
takes 128 channels and emits 64. Both use \(3\times3\) kernels with
padding that maintains \(7\times7\) resolution. Hidden state starts at
zero each trial; gate and candidate biases start at zero. Zero bias does
not force every gate to \(0.5\) on arbitrary input, because the
convolution weights still contribute. {[}D1{]}

A \textbf{write gate} close to zero carries the old hidden value
forward; close to one replaces it with the candidate. A \textbf{reset
gate} close to zero removes that old hidden component from the
candidate's input. It does not directly erase the carried state: the
carry term is governed by the write gate.

For example, let one hidden component be \(0.8\) and its candidate be
\(-0.2\). With \(w=0.1\), the new value is \(0.7\). With \(w=0.9\), it
is \(-0.1\). The same candidate can mean a minor revision or almost
complete replacement depending on the write gate. This example isolates
the final mixing operation; real candidates and gates depend on
neighboring features.

\subsection{5.3 What recurrence can do with
space}\label{what-recurrence-can-do-with-space}

A \(3\times3\) convolution gives each new hidden location access to a
local neighborhood of earlier hidden values and current inputs. Repeated
steps can propagate information across more of the hidden map. Through
the reset computation, the candidate also has indirect dependencies
beyond a single simple neighborhood, so ``one cell per frame'' would be
an oversimplified propagation law.

A network could learn to maintain an orientation representation at a
location, bring a later probe representation into contact with it, or
integrate evidence while preserving which region supplied it. These are
possible computations, not guaranteed emergent behaviors. Kernel sharing
favors reuse of local operations across positions. Padding, upstream
features and the final position-specific dense readout mean that the
whole model is not guaranteed to be perfectly translation invariant.

Starting from zero, the stated ConvGRU update keeps each hidden
component within \([-1,1]\) in exact arithmetic: the candidate is
bounded by tanh and each update is a convex mixture with the earlier
bounded component. This value bound does not establish nonexpansive
gradients; the gates and candidate themselves depend on hidden state. It
also does not directly apply to KDA matrices, whose values are signed
and unconstrained.

\subsection{5.4 KDA and ConvGRU store different
things}\label{kda-and-convgru-store-different-things}

{\def\LTcaptype{none} % do not increment counter
\begin{longtable}[]{@{}
  >{\raggedright\arraybackslash}p{(\linewidth - 4\tabcolsep) * \real{0.3333}}
  >{\raggedright\arraybackslash}p{(\linewidth - 4\tabcolsep) * \real{0.3333}}
  >{\raggedright\arraybackslash}p{(\linewidth - 4\tabcolsep) * \real{0.3333}}@{}}
\toprule\noalign{}
\begin{minipage}[b]{\linewidth}\raggedright
Question
\end{minipage} & \begin{minipage}[b]{\linewidth}\raggedright
Spatial KDA
\end{minipage} & \begin{minipage}[b]{\linewidth}\raggedright
Planned final ConvGRU
\end{minipage} \\
\midrule\noalign{}
\endhead
\bottomrule\noalign{}
\endlastfoot
Persistent state & Matrix per site and head & Hidden channel vector per
site \\
Write rule & Correct a key-addressed prediction error & Mix old state
with a nonlinear candidate \\
Retention control & Decay along key rows & Component-wise carry/write
mixing \\
Read rule & Query the matrix, then project & Expose hidden map to final
readout \\
Direct recurrent spatial mixing & No neighbor-state term in local matrix
update & Neighboring hidden values enter convolutions \\
Proposed role here & Memory inside three encoder scales & Spatial
recurrence after the deepest map \\
\end{longtable}
}

The new branch is a \textbf{KDA-plus-ConvGRU model}, not a contest in
which ConvGRU replaces KDA. The two forms of state may learn
complementary functions, redundant functions, or functions that
interfere. Their coexistence does not establish a clean biological
division into sensory memory and working memory.

\subsection{5.5 Three architectures that must not be
confused}\label{three-architectures-that-must-not-be-confused}

The \textbf{existing joint model} has KDA at three scales and a final
vector GRU. The \textbf{historical accumulator comparison} could put a
32-channel ConvGRU at those three scales, in place of KDA, while
retaining a final vector GRU. The \textbf{new planned branch} keeps the
KDA hierarchy and adds a 64-channel final ConvGRU in place of the
vector-recurrent route. {[}C1-C2; D1{]}

A performance result from the historical ConvGRU accumulator arm does
not measure the new branch. It differs in component placement, state
dimensions, training history and task exposure. Its existence helps
explain the design vocabulary, but it cannot stand in for the pending
implementation or its future results.

\subsection{5.6 What the new experiment could teach
us}\label{what-the-new-experiment-could-teach-us}

The proposed final hidden field has \(64\times7\times7=3,136\) scalars
per example, compared with 256 in the old final GRU. This is 12.25 times
as many final recurrent scalars; it is not a 12.25-fold increase in the
whole model's state, because the 215,808 KDA scalars remain. The
proposed combined count is 218,944, excluding buffers and training
intermediates.

The design warm-starts compatible CNN, KDA and task-head parameters from
a recorded global-GRU checkpoint, while initializing the new spatial
projection, ConvGRU and final projection. All learned parameters remain
trainable. Compatible optimizer state is to be migrated by parameter
name and shape. This is a planned training design, not a completion
receipt. {[}D1{]}

Preserving the task heads does not preserve their input representation:
their 256-dimensional inputs will now come from a newly initialized
route. Some initial disruption is therefore plausible. Later improvement
would show that this new branch can learn useful behavior under its
exposure. Without a matched continuing control, improvement would not
isolate the benefit of spatial recurrence from extra training, changed
capacity, initialization or other changes in the readout.

\newpage

\section{6. The seven two-frame sensory
tasks}\label{the-seven-two-frame-sensory-tasks}

\subsection{6.1 What this group
isolates}\label{what-this-group-isolates}

These tasks ask whether the system can extract and compare visual
information before we demand extended retention or complex spatial
instructions. Each trial contains exactly two presented images and one
final supervised answer. There is one primary condition, \texttt{mixed},
per task; difficulty values are mixed within that condition. Different
difficulty values are reporting strata, not additional primary cells.
{[}C4-C6{]}

All seven tasks still use the recurrent architecture. However, the final
three-frame stack contains both observed images. Good performance here
does not establish long-term retention. These tasks are useful
prerequisites and controls: weak elementary motion makes a failure on
motion-duration accumulation hard to attribute specifically to memory.

A binary label is not universally ``change.'' Five sensory tasks ask
\textbf{which interval} contains a particular property. Signed
orientation asks \textbf{which sign} the rotation has. Motion asks
\textbf{which direction}. These meanings must remain attached to the
corresponding task head.

\subsection{6.2 Motion direction: correspondence across two
images}\label{motion-direction-correspondence-across-two-images}

\textbf{Task key:} \texttt{motion\_direction}. \textbf{Labels:} right
\(=0\), up \(=1\), left \(=2\), down \(=3\). \textbf{Chance:} 25\% for
balanced four-way decisions.

Two random-dot images depict one of four cardinal displacement
directions. There are 256 dot identities in a periodic square domain; a
circular aperture determines which rendered dots are visible. Half the
identities survive between images and move in the chosen direction. The
other half are regenerated independently. The displacement is one, two
or three pixels. Thus ``256 dots'' does not mean all 256 are
simultaneously visible inside the aperture. {[}C5{]}

\textbf{Example.} If surviving dots move two pixels to the left, the
answer is class 2. Randomly replaced dots provide distractor
correspondences. A model needs a statistic of ordered spatial
correspondence, not the absolute location of one bright dot. Averaging
the two frames symmetrically would discard the sign of motion.

The useful invariant is the displacement direction despite random
initial positions and replacements. The renderer gives no intended
direction cue in the distribution of an individual frame. Nevertheless,
verifying that a trained model relies on the pair requires an actual
single-frame or order-control evaluation; it does not follow just from
naming the task.

\textbf{What success teaches us.} The model can read short-range motion
information in this raster regime. It does not establish accumulation
over eight transitions, selection among four patches, or change
detection between noncardinal directions. Those are distinct later
tasks.

\subsection{6.3 Signed orientation: compare axes, not pixel
phases}\label{signed-orientation-compare-axes-not-pixel-phases}

\textbf{Task key:} \texttt{orientation}. \textbf{Labels:} negative
axial-angle change \(=0\), positive change \(=1\).

A localized grating appears in each frame. The base orientation varies,
and the second orientation differs by \(\pm4\), \(\pm10\) or \(\pm22\)
degrees. The phases of the two gratings vary independently. Phase
changes move light and dark stripes without changing their orientation,
so simple pixel subtraction is not a sufficient orientation
representation. {[}C5{]}

Grating orientation is \emph{axial}: rotating by \(180^\circ\) gives the
same axis. A convenient signed comparison is

\[\Delta\theta=\tfrac12\operatorname{atan2}\!\left(\sin2(\theta_2-\theta_1),\cos2(\theta_2-\theta_1)\right).\]

For example, \(175^\circ\) followed by \(5^\circ\) can represent a
positive \(10^\circ\) axial change rather than a negative \(170^\circ\)
change. This formula explains the geometry; the model is not given the
numerical angles or this comparator.

The repository's older prose uses inconsistent
clockwise/counterclockwise wording. Use the renderer's numerical sign
rather than redefining labels from screen-coordinate intuition.
\textbf{Every trial has a nonzero change.} The task is not change versus
no change. High accuracy would establish signed comparison at these
sampled magnitudes, not a measured orientation threshold below four
degrees.

\subsection{6.4 Contrast: identify the stronger
modulation}\label{contrast-identify-the-stronger-modulation}

\textbf{Task key:} \texttt{contrast}. \textbf{Labels:} first frame
\(=0\), second frame \(=1\).

The two frames share a grating pattern. One has a baseline contrast, or
pedestal, of \(0.08\), \(0.18\) or \(0.30\); the other adds \(0.025\),
\(0.06\) or \(0.13\). The correct interval is the one with the larger
modulation around the background level. This is not simply which frame
is brighter on average. {[}C5{]}

\textbf{Example.} Frame 1 has contrast \(0.18\) and frame 2 has
\(0.24\). The increment is \(0.06\) and the correct label is 1. The same
increment from a pedestal of \(0.30\) creates a smaller relative change.
Report both pedestal and increment so a high mixed score does not
conceal weakness at subtle relative differences.

A plausible computation estimates orientation-insensitive modulation
energy and compares it across intervals. The network is not required to
implement that particular estimator. Strong contrast discrimination
demonstrates access to this image statistic; it does not by itself
demonstrate spatial attention or object recognition.

This task also illustrates a limitation of architectural intuition.
Normalization can alter contrast information, but one cannot infer
complete contrast removal from the presence of GroupNorm: convolution
biases, nonlinearities, spatial structure and multiple channels can
retain usable cues. The learned system's behavior settles whether the
available route is sufficient.

\subsection{6.5 Spatial frequency: compare stripe
density}\label{spatial-frequency-compare-stripe-density}

\textbf{Task key:} \texttt{spatial\_frequency}. \textbf{Labels:} first
frame \(=0\), second frame \(=1\).

The higher-frequency grating has more cycles over the same image extent.
The baseline is sampled between four and nine cycles per image, and the
higher interval multiplies it by \(2^\delta\), where \(\delta\) is
\(0.08\), \(0.18\) or \(0.35\) octaves. Independent phases prevent the
decision from reducing to a fixed stripe alignment. {[}C5{]}

\textbf{Example.} With a six-cycle baseline and an increment of \(0.18\)
octaves, the higher frequency is approximately \(6\times2^{0.18}=6.80\)
cycles per image. If it appears first, the label is 0. An octave is a
frequency ratio, so the absolute difference depends on the baseline.

The network must distinguish orientation, phase and frequency rather
than confusing any pattern difference with the requested property.
Convolutional channels with different effective frequency sensitivity
could support the comparison, but the code contains no hand-built
frequency classifier.

Success establishes discrimination in this sampled range and
rasterization. It is not scene recognition, and it is not proof of
human-like frequency tuning. A useful performance report separates all
three increments rather than treating the mixed average as one universal
frequency threshold.

\subsection{6.6 Chromatic increment: a specified colour
direction}\label{chromatic-increment-a-specified-colour-direction}

\textbf{Task key:} \texttt{chromatic\_increment}. \textbf{Labels:} first
frame \(=0\), second frame \(=1\).

A soft-edged patch appears on gray. One interval adds an increment of
\(0.018\), \(0.045\) or \(0.1\) along a specified linear-RGB colour
axis. Let \(w=(0.2126,0.7152,0.0722)\) be the renderer's luminance
weights. Its colour direction is proportional to

\[a=(0.7152,-0.2126,0),\qquad w^\top a=0.\]

After normalizing \(a\), the target colour is \(c+\delta a\). Under this
numerical luminance definition, the increment changes colour without
changing the patch's nominal luminance. Small image noise remains. This
is a renderer-defined colour discrimination, not a calibrated claim
about a particular display or human isoluminance. {[}C5{]}

\textbf{Example.} If the second patch has the positive increment, its
label is 1 even if a viewer informally describes the difference as
``redder'' or ``less green.'' The label is the specified axis increment,
not an arbitrary colour preference.

Because the baseline RGB colour varies, learning a fixed absolute target
colour is insufficient as a general solution. The model should compare
the intervals along the relevant direction. Strong performance supports
sensitivity to the defined colour contrast; it does not imply
categorical colour naming, constancy under illumination changes, or
spatial selection.

\subsection{6.7 Contour grouping: the arrangement is the
signal}\label{contour-grouping-the-arrangement-is-the-signal}

\textbf{Task key:} \texttt{contour}. \textbf{Labels:} first frame
\(=0\), second frame \(=1\).

Both images contain 32 small Gabor elements. In the structured interval,
seven are arranged and oriented along a curved path. Their alignment
jitter is two, eight or sixteen degrees. The control interval permutes
the orientation assignments across the same positions. The two images
therefore share the orientation multiset, but differ in which
orientations occur at which locations. {[}C5{]}

\textbf{Example.} One image contains seven appropriately aligned
elements along an arc; the other redistributes those orientations among
the 32 positions. If the aligned arc is in the first image, the answer
is 0. A histogram of orientation counts alone cannot distinguish the
intended manipulation because the multiset is matched.

The challenge is spatial organization: local orientation must be related
to neighboring positions and the path geometry. A network can exploit
local alignment statistics without developing a human-like percept of a
completed object boundary. That distinction is useful when interpreting
success.

Larger jitter weakens path alignment. Reporting only an overall accuracy
can hide whether the system handles the highly aligned case but fails
when the same path becomes less regular. A low score might reflect
insufficient spatial grouping, inadequate receptive-field use, difficult
optimization or weak local features; it is not directly a working-memory
capacity result.

\subsection{6.8 Natural spectral detail: a photograph task without
identity
memory}\label{natural-spectral-detail-a-photograph-task-without-identity-memory}

\textbf{Task key:} \texttt{natural\_spectrum}. \textbf{Labels:} first
frame \(=0\), second frame \(=1\).

Both intervals derive from the same natural photograph crop. The
renderer converts to luminance, preserves Fourier phase, and changes the
relative frequency content. Its multiplier is proportional to
\((r/0.1)^{-\beta}\) at radial frequency \(r\), with DC treated
separately. A smaller \(\beta\) retains more high-frequency detail
relative to low-frequency structure. The two beta values differ by
\(0.15\), \(0.30\) or \(0.60\). The resulting images share mean and RMS
contrast. {[}C6{]}

\textbf{Example.} If the first interval uses the smaller beta, it is the
correct answer even though both images show the same scene. ``More
detail'' here has an exact spectral meaning, rather than a judgment
about semantic informativeness.

Matching mean and RMS prevents the intended discrimination from being a
trivial comparison of overall brightness or contrast energy. Phase
preservation maintains much of the scene's spatial structure while
redistributing spectral energy. The model can succeed by reading texture
or edge-scale statistics; it need not name the scene or remember its
identity across a list.

The BSDS500 photograph identities are split into train, validation and
test pools. These same identity boundaries are shared with scene
recognition. Generalization to held-out photographs is therefore
different from sampling another random crop of a photograph already used
in training. This task and recognition should never be merged into one
vague ``natural-image ability.''

\newpage

\section{7. Six tasks with spatial instructions or longer
sequences}\label{six-tasks-with-spatial-instructions-or-longer-sequences}

\begin{figure}
\centering
\includegraphics[width=1\linewidth,height=\textheight,keepaspectratio,alt={The six longer task formats differ in cue timing, evidence, and report. D denotes blanks, B baseline motion, N study load, and H probe repetitions.}]{tmp/pdfs/trial_timelines.pdf}
\caption{The six longer task formats differ in cue timing, evidence, and
report. D denotes blanks, B baseline motion, N study load, and H probe
repetitions.}
\end{figure}

\subsection{7.1 Ring-cued orientation: location plus signed
change}\label{ring-cued-orientation-location-plus-signed-change}

\textbf{Task key:} \texttt{orientation\_ring}. \textbf{Labels:} negative
target rotation \(=0\), positive target rotation \(=1\). \textbf{Primary
cell:} \texttt{D0}.

The trial is \textbf{one cue frame, two sample frames, one probe frame}.
Four Gabor patches occupy the established four locations. A ring marks
the target during the cue and sample frames. Each patch rotates between
sample and probe; the foils have independently chosen signs. Rotation
magnitudes are 15, 30 or 45 degrees. There are no inserted retention
blanks in this suite cell. {[}C7{]}

\textbf{Example.} The ring selects the upper-right patch. That patch
rotates by \(+30^\circ\), while two foils rotate negatively and another
rotates positively. The answer is 1 because the target's sign is
positive. A majority vote over all patch signs is not the rule.

This combines location selection with signed orientation comparison. It
does not ask the model to interpret a plus/minus instruction, and it
contains no unchanged target negatives. Those distinctions make it
different from \texttt{orientation\_cued}, despite both having binary
outputs.

The final input stack includes the two samples and probe together. A
good solution can therefore be largely a spatially selective image
comparison at the final step. This is why historical ring-first training
could serve as an acquisition rung, but ring success alone cannot
establish long-delay memory. In the current joint suite, it is one task
among thirteen; its inclusion does not impose a ring-first curriculum.

\subsection{7.2 Signed cued orientation: apply a rule to the chosen
location}\label{signed-cued-orientation-apply-a-rule-to-the-chosen-location}

\textbf{Task key:} \texttt{orientation\_cued}. \textbf{Labels:} target
rotation agrees with instruction \(=1\); unchanged or opposite-sign
target rotation \(=0\). \textbf{Cells:} \texttt{D0}, \texttt{D4},
\texttt{D12}, \texttt{D24}.

The trial is \textbf{cue 1, sample 2, blank delay \(D\), probe 1},
totaling \(D+4\) frames. A plus/minus glyph near one location
simultaneously identifies the target and the sign to report. It is
visible in the cue and both samples, then absent from the probe. The
four centers are \((27,27)\), \((73,27)\), \((27,73)\) and \((73,73)\)
in image coordinates. {[}C8{]}

Let \(s\in\{-1,+1\}\) be the glyph's sign and \(\Delta\theta_j\) the
target's signed rotation. The label rule is

\[y=\mathbf1[s\Delta\theta_j>0].\]

\textbf{Examples.} A minus instruction with a target rotation of
\(-15^\circ\) is positive. The same instruction with \(+15^\circ\) is
negative. A zero rotation is also negative. Thus a real physical target
change can correctly receive a negative report.

The renderer constructs a sign-relative rotation multiset before
assigning the target, preserving aligned, opposite and unchanged
locations on every trial. This makes a global inventory of change signs
insufficient: the network needs the relation among instruction, target
location and target rotation.

A successful computation could retain target orientation and the
requested sign until the probe. It could instead keep a richer scene and
select later. The architecture does not stipulate the strategy. Failures
can arise from glyph recognition, location association, orientation
precision, retention or final rule application. Looking at all four
delays and comparing with the ring task helps organize these
possibilities, but does not uniquely diagnose them.

\subsection{7.3 Cued motion duration: count evidence across
transitions}\label{cued-motion-duration-count-evidence-across-transitions}

\textbf{Task key:} \texttt{motion\_duration\_cued}. \textbf{Labels:}
right/up/left/down \(=0/1/2/3\). \textbf{Cells:} the same four delays.

The trial is \textbf{cue 1, reference 1, moving frames 8, blanks \(D\),
report 1}, totaling \(D+11\) frames. Four dot patches have independently
generated direction schedules. A ring selects one patch and remains
visible during the reference and moving frames. Each patch has 32 dots;
on every transition, sixteen randomly chosen dots and any
boundary-crossing dots are replaced. Step size is \(0.8\), \(1.2\) or
\(1.6\) pixels. {[}C8{]}

The correct answer is the unique direction with the largest number of
the target's eight transitions:

\[y=\arg\max_{d\in\{0,1,2,3\}}\sum_{t=1}^{8}\mathbf1[d_t=d].\]

This formula defines the ground truth. The network receives images, not
the direction sequence or a counter.

\textbf{Example.} Consider right, up, right, left, down, right, up,
left. Counts are \((3,2,2,1)\), so the answer is right, even though the
last transition is left. The corresponding net displacement is not a
categorical count of the most frequent direction either. A strategy
based on final motion, one salient jump, or pooled evidence from all
patches can fail.

The task demands elementary motion estimation, target selection,
temporal accumulation and retention until report. The report frame
supplies no new motion evidence. Strong two-frame motion performance is
therefore a useful prerequisite but not a solution to this task.

A useful diagnostic separates count-winner margin from recency. A narrow
margin is intrinsically less tolerant of one misread transition; a
strong dependence on the last transition may indicate recency weighting.
Neither observation alone identifies whether the source is encoding,
recurrence or the final readout.

\subsection{7.4 Krauzlis-style motion change: reject a real foil
event}\label{krauzlis-style-motion-change-reject-a-real-foil-event}

\textbf{Task key:} \texttt{krauzlis\_cued\_motion}. \textbf{Labels:}
change at the cued patch \(=1\); foil-only change or catch \(=0\).
\textbf{Cells:} \texttt{B12}, \texttt{B20}, \texttt{B28}.

There are two patches, centered at \((20,50)\) and \((80,50)\). The
trial contains \textbf{cue 2, fixation-only 5, reference 1, baseline
transitions \(B\), postevent transitions 8, report 1}, totaling \(B+17\)
frames. The B values are baseline-motion lengths, not blank-memory
delays. The renderer defines a 100 Hz clock; its timing is a compressed
adaptation of Arcizet and Krauzlis's task, not a faithful recreation of
the entire experiment. {[}C8; P4{]}

The patches begin with mean directions separated by \(90^\circ\). An
event rotates one mean by \(\pm26^\circ\) or \(\pm28^\circ\), or
produces no change. Each patch has sixteen dots, subpixel bilinear
rendering, \(0.375\)-pixel steps, directional variability and finite dot
lifetimes. Motion is not restricted to four cardinal directions.

Over a complete 100-trial event cycle there are 57 target events, 29
foil events and 14 catches. \textbf{Example:} cue left, change right
means label 0, although a genuine motion change occurred. Cue left,
change left means label 1. No change means label 0 regardless of cue
side.

An always-positive model would achieve 57\% raw accuracy on that exact
mixture, 100\% target hits, and 100\% foil and catch false reports. Its
balanced accuracy would be 50\%. Hits alone would therefore make a
failed selective observer look excellent.

This task teaches the difference between detecting change and selecting
relevant change. Report target hits, foil false alarms and catch false
positives separately. A foil error may reflect poor spatial selection; a
catch error may reflect an overly liberal criterion or noisy change
estimation. Those are useful hypotheses, not automatic mechanistic
diagnoses. The model outputs one final decision, so it does not produce
a learned reaction-time distribution.

\subsection{7.5 Spatial binding: remember which feature belonged
where}\label{spatial-binding-remember-which-feature-belonged-where}

\textbf{Task key:} \texttt{spatial\_binding}. \textbf{Labels:} the
queried item participated in the exchange \(=1\); only foils exchanged
\(=0\). \textbf{Cells:} four delays.

The trial is \textbf{instruction 1, sample 2, blanks \(D\),
retrospective ring query 1, probe 1}, totaling \(D+5\) frames. No target
is marked during sample encoding. The four sample orientations come from
an inventory spaced \(45^\circ\) apart around an axial circle with a
random offset. At probe, exactly one pair of locations exchanges
orientations. The overall inventory stays the same. {[}C8{]}

\textbf{Example.} Suppose the four locations contain
\([10,55,100,145]^\circ\). After the delay, the query selects the first
location. Probe \([55,10,100,145]^\circ\) is positive because the
queried item exchanged with the second. Probe \([10,100,55,145]^\circ\)
is negative because only the second and third exchanged.

Both trials contain exactly two changed locations and exactly the same
set of orientation values as their sample. Neither an ``anything
changed'' detector nor an unordered bag of remembered orientations
solves the required distinction. The model needs information about the
association between orientation and location, accessible after the
retrospective query.

The late cue changes the informational requirement. The model cannot
know in advance which location will matter from an encoding precue. It
must preserve useful information across the initially eligible
locations, or adopt another strategy that still supports the eventual
query. Calling this ordinary precued selective storage would misdescribe
the trial.

Successful performance supports usable feature-location information for
these four familiar centers and sampled delays. It does not directly
estimate an unrestricted object capacity or generalization to new
spatial layouts. A query failure and a lost binding can both cause an
incorrect answer, so diagnostic probes must distinguish availability
from use.

\subsection{7.6 Scene recognition: exact membership after a
list}\label{scene-recognition-exact-membership-after-a-list}

\textbf{Task key:} \texttt{image\_recognition}. \textbf{Labels:} exact
study-list member \(=1\), nonmember \(=0\). \textbf{Cells:} every
combination of study load \(N\in\{0,4,12,24\}\) and probe hold
\(H\in\{3,4,5\}\), giving twelve cells.

The trial is \textbf{instruction 1, \(N\) distinct study images, blanks
3, \(H\) copies of one probe}, totaling \(N+4+H\) frames. Each study
image appears for one frame. Photographs are canonically cropped and
resized to \(100\times100\) RGB. Positive probes are byte-identical
copies of a study raster; negatives are distinct from all study rasters.
No instruction glyph overlays the photograph frames. {[}C8{]}

\textbf{Example.} After studying \([A,B,C,D]\), the repeated probe
\(C,C,C,C\) is positive. A repeated new image \(E,E,E,E\) is negative.
The repetition holds the same probe on screen; it does not create four
independent tests or four new studied items.

The third blank removes direct study-image access from the three-frame
stack before probe presentation. The model must retain some useful list
information elsewhere. It need not reconstruct every image. A learned
compressed familiarity or matching representation could suffice,
provided it generalizes to the held-out photograph pool and respects
membership.

For \(N=0\), every trial is negative. Report specificity and
false-positive rate; balanced accuracy and AUC are undefined because one
true class is missing. These three empty-list cells cannot raise a claim
about recognition of nonempty lists. For nonempty cells, compare load
and probe hold separately. A longer hold permits more recurrent
computation on the probe, but does not increase study exposure.

The task tests exact-raster membership. It does not require recognizing
the same scene from a new viewpoint, a different crop or altered
illumination. It is also not the spectral-detail task: recognition
compares the probe with a retained list, whereas spectral detail
compares two transformed views of the same source crop.

\subsection{7.7 The complete condition
map}\label{the-complete-condition-map}

{\def\LTcaptype{none} % do not increment counter
\begin{longtable}[]{@{}lrrr@{}}
\toprule\noalign{}
Task group & Tasks & Cells per task & Total cells \\
\midrule\noalign{}
\endhead
\bottomrule\noalign{}
\endlastfoot
Two-frame sensory & 7 & 1 & 7 \\
Ring orientation & 1 & 1 & 1 \\
Signed cued orientation & 1 & 4 & 4 \\
Cued motion duration & 1 & 4 & 4 \\
Krauzlis motion change & 1 & 3 & 3 \\
Spatial binding & 1 & 4 & 4 \\
Scene recognition & 1 & 12 & 12 \\
\textbf{Total} & \textbf{13} & & \textbf{35} \\
\end{longtable}
}

Thirty-two cells have both-class or multiclass scoring eligibility. The
three empty-list recognition cells remain visible but are excluded from
BA/AUC checkpoint selection. Difficulty strata, such as rotation
magnitude or motion displacement, sit within these cells rather than
changing the 35-cell total. {[}C4{]}

\newpage

\section{8. How the tasks train a shared recurrent
system}\label{how-the-tasks-train-a-shared-recurrent-system}

\subsection{8.1 One final loss, many earlier
computations}\label{one-final-loss-many-earlier-computations}

Each trial produces logits from the selected task head. Cross-entropy
penalizes low probability on the correct class:

\[L=-\frac1B\sum_{b=1}^{B}\log p_{b,y_b}.\]

Only the selected head participates in that task's computation. The
shared encoder, KDA modules and final recurrent route can receive
gradients from every task. ``All learned parameters are trainable''
means no component is deliberately frozen; it does not mean every
parameter receives a nonzero gradient on every batch. Inactive task
heads have no gradient for that update. {[}C3{]}

There are no supervised target-angle labels for the internal state, no
gate target saying ``retain now,'' and no extra output at every frame in
the deployed model. Full-sequence backpropagation propagates the final
decision's loss through the intervening state transitions. It can teach
an early gate to preserve information if doing so improves the eventual
answer. This is why task design and temporal gradient flow matter even
when the visible supervision is only a final class label.

For a generic recurrence \(z_t=f_\theta(z_{t-1},x_t)\), credit to an
early state involves products of state-transition Jacobians. Long paths
can attenuate or amplify gradients. Gating can provide useful carry
paths, but does not guarantee successful credit assignment. The local
KDA nonexpansion result in Chapter 2 is not a blanket proof about this
full training computation.

\subsection{8.2 Effective batch size and
microbatches}\label{effective-batch-size-and-microbatches}

The existing joint trainer uses effective batch size 32 through eight
microbatches of four. Each microbatch computes mean cross-entropy,
scales it by \(4/32\), and accumulates gradients before one Adam step.
All eight microbatches belong to the same chosen task and condition.
{[}C3{]}

This reduces peak batch memory while preserving the intended average
gradient, up to numerical differences, for this model's per-example
computation. It does not truncate a sequence: each microbatch
backpropagates through its complete trial. Microbatching across examples
and truncating through time are different operations.

The recorded optimizer is Adam with learning rate \(10^{-4}\), betas
\((0.9,0.999)\), epsilon \(10^{-8}\) and zero weight decay. There is one
parameter group, fp32 computation, finite-value checks, and no gradient
clipping. These describe this training recipe, not a claim that it is
universally optimal. The new branch's brief preserves these settings
while changing the final architecture. {[}C3; D1{]}

\subsection{8.3 Equal tasks do not mean equal conditions or
frames}\label{equal-tasks-do-not-mean-equal-conditions-or-frames}

A shuffled cycle contains one update from each of thirteen tasks. Within
a task, a shuffled queue cycles through its conditions. Thus task
exposure is balanced by optimizer updates, while a twelve-cell task
receives fewer updates per individual cell than a one-cell task over the
same total horizon. {[}C3{]}

Over 156 updates, there are twelve updates per task. A one-cell sensory
task receives all twelve in its single cell. A four-cell task receives
three per cell. Recognition receives one per cell. With 32 examples per
update, each task receives 384 episodes over this illustration, but its
allocation across conditions differs.

Episodes also differ in length. A two-frame task uses far fewer image
frames than a long motion or recognition trial. Equal episodes therefore
do not imply equal frame exposure or equal compute cost. When
interpreting learning speed, retain three counts: optimizer updates,
independent episodes, and processed frames.

Joint training creates a shared representational problem. A feature
useful for contrast might be irrelevant to binding; a gate useful for
sensory comparison may need a different pattern for retention. Loss
curves alone cannot establish destructive gradient interference. That
hypothesis needs a comparison or diagnostic designed to measure it.

\subsection{8.4 Acquisition, validation and
testing}\label{acquisition-validation-and-testing}

Training changes weights. Validation selects among checkpoints or checks
acquisition under a declared policy. Final testing evaluates the
selected result on a separate split. Repeatedly inspecting a fixed
validation set does not make each look independent new evidence.

The suite maintains separate task/condition streams. Sampling one
condition does not advance another. Its ordinary delay cells are
independent draws, not the same stimulus rendered at multiple delays. To
estimate the causal effect of delay on a matched scene, a separate
paired evaluation must hold the base stimulus fixed and change only the
delay. {[}C4{]}

The existing project's selection policy changed between training phases.
The first short run used a worst-task chance-normalized BA criterion;
later continuation used equal-task mean AUC with a BA tie-break. Those
are different optimization preferences. A minimum rewards the weakest
task but can be sensitive to a noisy floor; a mean can improve while
difficult tasks remain at chance. Neither criterion substitutes for
inspecting the complete task vector. {[}R1; C3{]}

The supplied note's step-715 test and step-2314 validation are
historical snapshots with different evaluation roles. They should not be
narrated as a paired architectural comparison or as the latest model
state. This book uses no new performance measurements.

\subsection{8.5 Accuracy, balanced accuracy and AUC answer different
questions}\label{accuracy-balanced-accuracy-and-auc-answer-different-questions}

Raw accuracy counts correct decisions. Balanced accuracy averages recall
across the true classes:

\[\mathrm{BA}=\frac1K\sum_{c=1}^{K}\frac{\mathrm{correct\ predictions\ in\ class}\ c}{\mathrm{examples\ in\ class}\ c}.\]

For binary tasks with both classes represented, a constant-class
predictor has BA \(0.5\), even when the class frequencies are unequal.
For the four-direction tasks, balanced chance is \(0.25\). A useful
cross-task scaling is

\[\mathrm{BA}_{\mathrm{normalized}}=\frac{\mathrm{BA}-1/K}{1-1/K}.\]

This maps chance to zero and perfect performance to one. It does not
equalize uncertainty or task difficulty. {[}C3{]}

Binary AUC measures ranking: the probability that a randomly chosen
positive receives a higher positive score than a randomly chosen
negative, with a half-credit convention for ties. Consider negatives
scored \(0.60,0.61\) and positives \(0.70,0.71\). At a threshold of
\(0.5\), every example is called positive, so BA is \(0.5\). Yet AUC is
1 because the ordering is perfect. This constructed example explains how
a ranking signal can coexist with a useless deployed threshold.

The project's multiclass AUC averages one-versus-rest AUCs. Its chance
reference is \(0.5\), not \(0.25\). AUC is not accuracy, and finding a
promising ranking does not itself demonstrate a corrected decision rule.
Threshold fitting must have its own training/validation separation.

\subsection{8.6 What should be in a learning
report}\label{what-should-be-in-a-learning-report}

For each task, keep the primary conditions, denominators, confusion
matrices and relevant difficulty strata visible. Preserve the special
Krauzlis event decomposition and empty-list recognition specificity.
Average eligible conditions within a task before computing an equal-task
summary, so recognition does not dominate merely because it has twelve
cells.

A trained seed is one trained realization. More test trials reduce
uncertainty about that realization's behavior on the stimulus
distribution; they do not replace independent training seeds.
Photograph-based uncertainty also needs to account for repeated use of a
finite source pool. A visually impressive average is useful only when
the underlying unit of evidence is clear.

\newpage

\section{9. From successful behavior to a
mechanism}\label{from-successful-behavior-to-a-mechanism}

\subsection{9.1 Availability and use are different
questions}\label{availability-and-use-are-different-questions}

Suppose a frozen-state decoder predicts sample orientation accurately
just before the probe, but the model's final answer is poor. This would
show that the extracted state contains decodable information for that
decoder and evaluation distribution. It would motivate investigating
access, comparison or report generation. It would not prove that the
deployed readout can use that information in the same way.

Conversely, a failed linear probe does not prove erasure. The code may
be nonlinear, distributed across states, or inaccessible to the chosen
analysis. Probe capacity, training data, target definition and
evaluation split all matter. For axial orientation, a decoder should
respect the doubled-angle geometry rather than treating \(1^\circ\) and
\(179^\circ\) as maximally distant.

A useful diagnosis follows the computation: test sensory representation,
target/instruction use, retained content, comparison and output. This
order prevents the word ``memory'' from becoming a catch-all explanation
for every sequence-task failure.

\subsection{9.2 Three senses of
attention}\label{three-senses-of-attention}

\textbf{Associative attention} is query-dependent retrieval from a
matrix. \textbf{Spatial selection} is behavior that depends on the
instructed location rather than irrelevant alternatives.
\textbf{Competitive spatial allocation} is a mechanism that explicitly
normalizes influence across locations, such as a spatial softmax.

Our KDA implements the first and can support learning the second. It
does not contain the third as an explicit spatial softmax. Its
\(\alpha\) is a row-retention gate, not a normalized probability of
attending to an image location. Calling every \(\alpha\) map ``attention
allocation'' would collapse different mechanisms into one word.

The recurrent vision transformer studied by Morgan, Albanna and Herman
uses memory-guided spatial attention and a learned wait/declare-change
policy. Its cue-validity manipulations and attention-bias interventions
concern a different computational substrate from our local associative
matrices and fixed final report. Its findings motivate comparison, but
not an assumption of equivalent dynamics. {[}P5{]}

\subsection{9.3 Perturbing an emission versus perturbing
memory}\label{perturbing-an-emission-versus-perturbing-memory}

The original note proposes computational interventions; it does not
report them as completed on the joint learner. An \textbf{emission
pulse} adds a spatially localized vector to \(O_t\) after the KDA output
projection. This changes what downstream processing receives at that
frame without directly overwriting that module's saved matrix. It can
nevertheless have lasting effects through higher-scale KDA and final
recurrent state.

A \textbf{state pulse} instead changes the carried matrix. For a
rank-one perturbation \(\Delta S=\lambda ab^\top\), the immediate raw
read changes by

\[\Delta o=\Delta S^\top q=\lambda(a^\top q)b.\]

If \(a\) is orthogonal to \(q\), the immediate raw read is unchanged
even though memory was altered. A later query may expose the change.
This is why perturbation size alone is not an adequate manipulation
check.

These are interventions on mathematical representations. Their dose has
no established conversion to electrode current, recruited neurons or
tissue extent. Signed KDA coordinates also make ``add a positive number
to every channel'' an ambiguous analogue of physiological excitation. A
direction should be calibrated and its achieved effect measured.
{[}R1{]}

\subsection{9.4 What a controlled causal comparison
needs}\label{what-a-controlled-causal-comparison-needs}

Compare identical episodes under sham and perturbation, with a frozen
checkpoint and a verified analysis wrapper. Specify the altered
substrate, scale, location, temporal phase and dose. Distinguish sample
encoding from the first two nominal blanks, which retain sample pixels
in the raw stack, and from later genuinely sample-free processing.

Compare target, foil and off-target sites, and match the intervention
geometry and magnitude as far as the representation allows. Include
patterns that test whether an effect is specific to the chosen feature
direction rather than generic disruption. This is an experimental design
lesson, not authorization to run an intervention.

Behavioral interpretation must respect the label. Increasing reports of
an uncued change is a false-alarm increase in the Krauzlis task. An
opposite-sign orientation change is negative under an incompatible sign
instruction. More positive reports are not automatically better
perception.

For a binary task, separate hits and false alarms. Where a
signal-detection model is appropriate,

\[d'=\Phi^{-1}(H)-\Phi^{-1}(F),\qquad
c=-\tfrac12[\Phi^{-1}(H)+\Phi^{-1}(F)].\]

Sensitivity \(d'\) and criterion \(c\) distinguish discrimination from
reporting tendency under the model's assumptions. Rates of zero or one
require a declared finite-sample correction. Four-choice direction
judgments should primarily use multiclass confusion, not an unexplained
binary conversion.

A null perturbation effect also needs interpretation. The manipulation
may have missed the relevant representation, been hidden from the
current query, or been compensated by another route. Necessity,
sufficiency, decodability and normal use are different claims.
Well-designed tests make those distinctions visible rather than turning
every outcome into confirmation.

\newpage

\section{10. Exercises with worked
answers}\label{exercises-with-worked-answers}

\subsection{10.1 Work the associative
update}\label{work-the-associative-update}

\textbf{Question.} Begin with zero state, a unit key \(k=(1,0)^\top\),
value \(v=(3,2)^\top\), query equal to the key, retention one and write
strength \(0.5\). Present the same key/value twice. What are the two
reads?

\textbf{Answer.} The first prediction is zero. The first write makes the
first row \((1.5,1)\), which is also the first read. At the second step,
the error is \((1.5,1)\) and half of it is added. The second read is
\((2.25,1.5)\). Repeated writes approach \((3,2)\) rather than adding it
without bound. Retention one is an idealized limiting case used to
isolate the delta correction.

\subsection{10.2 Find the interference}\label{find-the-interference}

\textbf{Question.} A write has \(\beta=0.5\), unit key \(k=(1,0)^\top\)
and error \(e=(2,-4)^\top\). What changes for query \(q'=(0,1)^\top\)?
What about \(q''=(1,1)^\top/\sqrt2\)?

\textbf{Answer.} The orthogonal query changes by zero. The diagonal
query changes by \(0.5(1/\sqrt2)(2,-4)^\top=(1/\sqrt2,-\sqrt2)^\top\).
The same write has different effects depending on address overlap; no
spatial change is needed for interference within a site's associative
code.

\subsection{10.3 Read a gate correctly}\label{read-a-gate-correctly}

\textbf{Question.} A proposed ConvGRU cell has write gate near zero and
reset gate near zero. Does reset erase its old state?

\textbf{Answer.} No.~Reset removes old-state input to the candidate
computation, but the near-zero write gate preserves the old state in the
final mixture. If the write gate were near one, the state would instead
be replaced by the largely current-input-driven candidate.

\subsection{10.4 Identify the changed architectural
hypothesis}\label{identify-the-changed-architectural-hypothesis}

\textbf{Question.} The new branch improves binding. Can we conclude that
replacing KDA with ConvGRU improved binding?

\textbf{Answer.} No.~The new branch keeps KDA and changes the final
recurrent readout. Moreover, its warm start and additional exposure mean
that improvement without a matched control would not isolate the
architecture's causal contribution. The defensible observation would be
that the new combined branch acquired better binding under its measured
training history.

\subsection{10.5 Interpret a signed cue}\label{interpret-a-signed-cue}

\textbf{Question.} A minus glyph selects the lower-left patch. Its
rotation is positive, while an uncued patch rotates negatively. Is the
correct report positive?

\textbf{Answer.} No.~The rule applies to the selected patch. Its
rotation disagrees with the requested sign, so the label is zero.
Selecting a foil with the desired rotation is an error even though the
desired motion exists somewhere in the image.

\subsection{10.6 Separate binding from global
change}\label{separate-binding-from-global-change}

\textbf{Question.} In a binding negative trial, did nothing change?

\textbf{Answer.} Two foil locations exchanged orientations. The queried
location was preserved. Every trial contains a swap, and the orientation
inventory is unchanged. Global change detection or an unordered
orientation inventory cannot provide the required answer.

\subsection{10.7 Audit an apparently impressive
score}\label{audit-an-apparently-impressive-score}

\textbf{Question.} A Krauzlis model reports ``change'' on every trial
and is described as achieving perfect target detection. What is missing?

\textbf{Answer.} Its foil and catch false-positive rates are also 100\%.
On the prescribed event mixture, raw accuracy is 57\% and BA is 50\%.
Perfect hit rate here is compatible with complete failure to
discriminate or select.

\subsection{10.8 Count evidence, not repeated
frames}\label{count-evidence-not-repeated-frames}

\textbf{Question.} A recognition trial has \(N=12\), \(H=4\). How many
frames and independent episodes does it contain? What changes if \(H\)
becomes five?

\textbf{Answer.} It has \(12+4+4=20\) frames and one episode. Increasing
\(H\) to five creates 21 frames in the episode, not another independent
membership judgment. It permits an additional recurrent update on the
same probe.

\subsection{10.9 Distinguish retained information from successful
behavior}\label{distinguish-retained-information-from-successful-behavior}

\textbf{Question.} A strong pre-probe decoder reads the sample
orientation from KDA, but the deployed model answers at chance. What
should the next explanation focus on?

\textbf{Answer.} The result makes complete absence of decodable sample
information less plausible for the tested state and decoder. It
motivates testing whether the normal query, final recurrent route and
head access and use that information, and whether the correct location
and instruction are represented. It does not by itself prove which of
those stages failed.

\newpage

\section{Reference sheet and source
map}\label{reference-sheet-and-source-map}

\subsection{Notation at a glance}\label{notation-at-a-glance}

{\def\LTcaptype{none} % do not increment counter
\begin{longtable}[]{@{}ll@{}}
\toprule\noalign{}
Symbol & Meaning in this book \\
\midrule\noalign{}
\endhead
\bottomrule\noalign{}
\endlastfoot
\(B,T\) & Batch size and presented sequence length \\
\(D\) & Inserted blank frames in delay tasks \\
\(B\) in \path & B12/B20/B28 \\
\(N,H\) in recognition & Study-list load and repeated-probe hold \\
\(S_t\) & KDA key-by-value matrix at a site/head \\
\(q_t,k_t,v_t\) & Query, key and incoming value \\
\(\alpha_t,\beta_t\) & Row retention and scalar write strength \\
\(O_t,F_t\) & KDA emission and deepest combined visual map \\
\(h_t,H_t\) & Vector GRU state and spatial ConvGRU state \\
\(w_t,r_t\) & ConvGRU candidate-write and reset gates \\
\(z_t\) & Old-state mixing gate in PyTorch's vector GRU \\
\end{longtable}
}

\subsection{Implementation and design
sources}\label{implementation-and-design-sources}

Paths below are relative to the supplied
\textbf{VisualAttentionWorkingMemory} repository. They are references
for readers, not instructions to execute the historical plans they
contain. A separate \texttt{source\_manifest.json} records SHA-256
identities of the local sources used for this edition.

\textbf{C1 - Associative update and historical accumulator cells.}
\path\textbar PreAttentiveVision/TemporalIntegration/accumulators.py\textbar:
\texttt{kda\_update}, \texttt{SpatialKDA}, \texttt{SpatialConvGRU}.
Defines the local equations, packed channel layout, normalization, gate
initialization and historical 32-channel ConvGRU. Chapters 2-3 derive
the code's algebra rather than relying on a language-model paper for
implementation-specific claims.

\textbf{C2 - Existing hierarchy and readout.}
\path\textbar WorkingMemory/PlainBaseline/accum.py\textbar:
\texttt{AccumulatorBaseline}, \texttt{frames}, \texttt{encode\_frame},
\texttt{forward}. Defines convolution widths, concatenation, centered
stacking, dense projection, vector GRU and selected head.

\textbf{C3 - Existing joint learning.}
\path\textbar SecondPass/JointTraining/worker.py\textbar{} and
\texttt{core.py\textbackslash{}path\textbar{}:\ model/optimizer\ construction,\ microbatch\ accumulation,\ scheduler,\ checkpoint\ state\ and\ scoring.\ The\ continuation\ policy\ is\ documented\ separately\ in\ \textbar{}continuation\_v3.py}
and \texttt{AMENDMENT\_V3.md}.

\textbf{C4 - Complete task inventory.}
\path\textbar SecondPass/TaskSuite/README.md\textbar,
\texttt{catalog.json} and \texttt{suite.py}. Defines all thirteen tasks,
35 cells, image-only interface, stream separation and scoring
eligibility. Historical authorization text is not an instruction to
begin a run.

\textbf{C5 - Synthetic two-frame renderers.}
\path\textbar PreAttentiveVision/neuroscience\_stimuli.py\textbar:
\texttt{CardinalMotionStream} and \texttt{TaskStream}. Defines cardinal
motion, signed orientation, contrast, spatial frequency, chromatic
increment and contour grouping.

\textbf{C6 - Photograph spectral task.}
\path\textbar PreAttentiveVision/natural\_stimuli.py\textbar:
\texttt{NaturalSpectrum}. Defines source partitions, crop generation and
frequency manipulation.

\textbf{C7 - Ring orientation.}
\path\textbar WorkingMemory/PlainBaseline/variants.py\textbar:
\texttt{VariantStream}. The suite uses its ring variant only; other
diagnostic rungs are not quietly added to the training task list.

\textbf{C8 - Spatial and sequence tasks.}
\path\textbar WorkingMemory/SpatialTaskBattery/stimuli.py\textbar:
\texttt{SpatialBatteryStream}, \texttt{RecognitionImages},
\texttt{frame\_count}. Defines signed cued orientation, motion duration,
Krauzlis change, binding and recognition. The motion schedule helper is
in \path\textbar WorkingMemory/stimuli.py\textbar.

\textbf{D1 - New branch, design only.}
\path\textbar SecondPass/SpatialReadout/BRIEF.md\textbar, read 24
September 2026, together with the user's clarification that another
agent is building the ConvGRU and it is not implemented yet. The reading
treats it as a specification and does not infer launch or success from
its authorization language.

\textbf{R1 - Supplied research note.} \emph{Spatial KDA memory:
architecture and causal tests}, 23 September 2026, supplied as
\texttt{87252611-f998-4054-8a87-a46d9a26e873.pdf}. Corresponding source
and static architecture audit are in
\path\textbar SecondPass/JointTraining/TechnicalReport/architecture\_microstimulation.md\textbar{}
and \texttt{architecture\_metrics\_notes.md}. Its measured scores retain
their original dates and evaluation roles; this textbook does not update
them.

\subsection{Primary literature and
documentation}\label{primary-literature-and-documentation}

\textbf{P1. Kimi Team (2025).}
\href{https://arxiv.org/html/2510.26692v1}{\emph{Kimi Linear: An
Expressive, Efficient Attention Architecture}}, pinned v1. Background
for the KDA name and fine-grained gated-delta lineage. The visual
implementation's details are grounded in C1.

\textbf{P2. Ballas, Yao, Pal and Courville (2016).}
\href{https://arxiv.org/abs/1511.06432}{\emph{Delving Deeper into
Convolutional Networks for Learning Video Representations}}. Primary
architectural precedent for convolutional GRU recurrence on video
features. The planned branch's exact specification is D1.

\textbf{P3. PyTorch documentation.}
\href{https://docs.pytorch.org/docs/main/generated/torch.nn.GRU.html}{\emph{GRU}},
accessed 24 September 2026. Source for the vector GRU equation and
reset-placement convention; this citation is not a claim that the
project uses the newest documentation's package version.

\textbf{P4. Arcizet and Krauzlis (2018).}
\href{https://journals.plos.org/plosbiology/article?id=10.1371/journal.pbio.2005930}{\emph{Covert
spatial selection in primate basal ganglia}}. Experimental provenance
for the local two-patch adaptation. C8, rather than the biological
paper, is authoritative for our rasterization, compressed timing and
label generation.

\textbf{P5. Morgan, Albanna and Herman (2025).}
\href{https://arxiv.org/html/2502.10955v1}{\emph{A recurrent vision
transformer shows signatures of primate visual attention}}, pinned v1.
Functional comparison for memory-guided spatial attention and causal
attention-bias tests. It does not establish the same mechanism in our
KDA learner.

\subsection{What this edition
establishes}\label{what-this-edition-establishes}

This edition explains the implemented computation and task definitions,
derives consequences of their equations, and teaches the pending
architectural proposal. It distinguishes illustrative arithmetic,
source-defined design and measured evidence throughout. Whether the
proposed ConvGRU improves any task remains a question for the separately
implemented and evaluated experiment.

\end{document}
